\documentclass[11pt,a4paper]{article}
\usepackage[utf8]{inputenc}
\usepackage[french]{babel}
\usepackage{amsmath,amssymb,amsthm}
\usepackage{geometry}
\geometry{hmargin=2.5cm,vmargin=3cm}
\usepackage{tikz}
\usetikzlibrary{cd,positioning,shapes}
\usepackage{listings}
\usepackage{xcolor}

\lstset{
    language=Python,
    basicstyle=\ttfamily\small,
    keywordstyle=\color{blue}\bfseries,
    stringstyle=\color{red},
    commentstyle=\color{gray}\itshape,
    numbers=left,
    numberstyle=\tiny\color{gray},
    breaklines=true,
    frame=single,
    showstringspaces=false,
    literate={é}{{\'e}}1 {è}{{\`e}}1 {à}{{\`a}}1 {ê}{{\^e}}1 {î}{{\^i}}1 {ç}{{\c{c}}}1 {ï}{{\"i}}1 {É}{{\'E}}1
}

\title{\Huge THÉORÈME DE CONVERGENCE DOMINÉE (TCD) \\ \large Cours magistral, exercices corrigés et travaux pratiques}
\author{\Large Charles EDOU NZE}
\date{\today}

\begin{document}
\maketitle
\tableofcontents
\newpage

\part{Cours Magistral}

\section{Jalon 69 : Théorème de convergence dominée (TCD)}

\section{Genèse et Intuition Physique}

\begin{figure}[h]
\centering
\begin{tikzpicture}
    % Axes
    \draw[->, thick] (-1,0) -- (6,0) node[right] {$x$};
    \draw[->, thick] (0,-2.5) -- (0,3) node[above] {$y$};

    % Fonction dominatrice g
    \draw[green!50!black, very thick, dashed, domain=0.5:5.5, samples=100] plot (\x, {2.5*exp(-0.5*(\x-3)*(\x-3)) + 0.2});
    \node[green!50!black] at (3, 3.1) {$g(x)$ (enveloppe int\'egrable)};

    % Fonction -g
    \draw[green!50!black, very thick, dashed, domain=0.5:5.5, samples=100] plot (\x, {-2.5*exp(-0.5*(\x-3)*(\x-3)) - 0.2});
    \node[green!50!black] at (3, -3.1) {$-g(x)$};

    % f_1
    \draw[blue, thick, domain=0.5:5.5, samples=100] plot (\x, {(2.5*exp(-0.5*(\x-3)*(\x-3)) + 0.2) * sin(\x*150)});
    \node[blue] at (5.5, 1.5) {$f_1$};

    % f_n
    \draw[red, thick, domain=0.5:5.5, samples=200] plot (\x, {(2.5*exp(-0.5*(\x-3)*(\x-3)) + 0.2) * sin(\x*400) * exp(-0.1*\x)});
    \node[red] at (5.5, -1) {$f_n$};

\end{tikzpicture}
\caption{Visualisation de la domination : Les fonctions $f_n$ oscillent fortement, mais restent enfermées entre $g$ et $-g$. L'aire sous $g$ étant finie, les intégrales des $f_n$ ne peuvent pas s'échapper.}
\end{figure}


Le théorème de convergence dominée (TCD) de Lebesgue est sans doute le résultat le plus emblématique et le plus utilisé de toute la théorie de la mesure. Il offre une réponse définitive et élégante au problème de l'interversion de la limite et de l'intégrale.
La théorie de Riemann exigeait une convergence uniforme, une condition souvent trop rigide, notamment en mécanique quantique ou en théorie des probabilités où des singularités isolées apparaissent. Lebesgue a démontré que si toute l'agitation des fonctions $f_n$ est contenue, en valeur absolue, sous une enveloppe mesurable et intégrable $g$, alors l'interversion est valide, indépendamment des variations microscopiques locales.

\section{Définitions, Théorèmes et Exemples}

\section{Théorème de Lebesgue (TCD)}

Soit $(X, \mathcal{F}, \mu)$ un espace mesuré.
Soit $(f_n)_{n \in \mathbb{N}}$ une suite de fonctions mesurables à valeurs réelles ou complexes.
Supposons que :
1. La suite $(f_n)$ converge presque partout vers une fonction $f$.
2. Il existe une fonction $g : X \to [0, +\infty]$ \textbf{intégrable} (i.e. $\int g d\mu < +\infty$) telle que :
   $$\forall n \in \mathbb{N}, \quad |f_n| \le g \quad \text{presque partout.}$$

Alors :
- $f$ est intégrable sur $X$.
- $\lim_{n \to \infty} \int_X f_n d\mu = \int_X f d\mu$.
- $\lim_{n \to \infty} \int_X |f_n - f| d\mu = 0$ (Convergence dans $L^1$).

\section{Exemples Concrets et Immédiats}

\textbf{Exemple 1 : L'enveloppe exponentielle}
Soit $f_n(x) = e^{-x} \cos(nx)$ sur $X = [0, +\infty[$.
On a $|f_n(x)| \le e^{-x} = g(x)$. La fonction $g(x) = e^{-x}$ est intégrable (intégrale valant 1).
Bien que $\cos(nx)$ n'ait pas de limite simple, cet exemple illustre la notion de domination.

\textbf{Exemple 2 : Calcul avec domination}
La limite simple de $f_n(x) = \frac{n\sin(x/n)}{x(1+x^2)}$ est $\frac{1}{1+x^2}$.
Domination : $|f_n(x)| \le \frac{x/n \cdot n}{x(1+x^2)} = \frac{1}{1+x^2}$.
La fonction $g(x) = \frac{1}{1+x^2}$ est intégrable, son intégrale est $\pi/2$.

\textbf{Exemple 3 : Défaut de domination (la bosse glissante)}
Soit $f_n(x) = n \mathbf{1}_{]0, 1/n[}(x)$ sur $]0, 1[$.
Limite simple : $\forall x \in ]0, 1[, \lim f_n(x) = 0$.
Mais $\int_0^1 f_n(x) dx = n \times \frac{1}{n} = 1$. L'intégrale de la limite est 0.
Pourquoi ? Car le supremum sur $n$ est $g(x) = \sup_n f_n(x) = \approx 1/x$, qui n'est pas intégrable en 0. L'hypothèse de domination n'est pas satisfaite.

\textbf{Exemple 4 : Fonctions puissances}
$f_n(x) = x^n$ sur $[0, 1[$.
Limite simple $f(x) = 0$.
Domination : $|f_n(x)| \le 1$. La constante 1 est intégrable sur le segment de longueur 1.
L'intégrale est $\lim \frac{1}{n+1} = 0 = \int 0 dx$.

\textbf{Exemple 5 : Paramètre de chaleur}
$f_n(x) = e^{-nx^2}$ sur $]0, 1[$.
Limite simple $0$. Domination : $|f_n(x)| \le e^{-x^2} \le 1$, intégrable sur $]0, 1[$.

\section{Démonstrations}

La démonstration s'appuie sur le lemme de Fatou, déjà étudié (Jalon 68).

1. \textbf{Intégrabilité de f :}
   Puisque $|f_n| \le g$ presque partout, en passant à la limite (la fonction valeur absolue étant continue), on obtient $|f| \le g$ presque partout. Par monotonie de l'intégrale et l'hypothèse $\int g d\mu < \infty$, on conclut que $\int |f| d\mu < \infty$. Donc $f$ est intégrable.

2. \textbf{Première application de Fatou (Minoration) :}
   Considérons la suite de fonctions $h_n = g + f_n$. Par l'hypothèse de domination, $f_n \ge -g$, donc $h_n \ge 0$.
   Appliquons le lemme de Fatou aux $(h_n)$ :
   $$\int \liminf_{n \to \infty} (g + f_n) d\mu \le \liminf_{n \to \infty} \int (g + f_n) d\mu$$
   Puisque $f_n \to f$ p.p., $\liminf (g + f_n) = g + f$.
   Ainsi, $\int (g + f) d\mu \le \int g d\mu + \liminf \int f_n d\mu$.
   Comme $g$ est intégrable, $\int g d\mu$ est un nombre fini, on peut le soustraire :
   $$\int f d\mu \le \liminf_{n \to \infty} \int f_n d\mu$$

3. \textbf{Seconde application de Fatou (Majoration) :}
   Considérons la suite $k_n = g - f_n$. Par domination, $f_n \le g$, donc $k_n \ge 0$.
   De même, avec Fatou :
   $$\int \liminf_{n \to \infty} (g - f_n) d\mu \le \liminf_{n \to \infty} \int (g - f_n) d\mu$$
   $$\int (g - f) d\mu \le \int g d\mu - \limsup_{n \to \infty} \int f_n d\mu$$
   En soustrayant à nouveau $\int g d\mu$ :
   $$-\int f d\mu \le -\limsup_{n \to \infty} \int f_n d\mu \implies \limsup_{n \to \infty} \int f_n d\mu \le \int f d\mu$$

4. \textbf{Conclusion :}
   Des étapes 2 et 3, on déduit l'encadrement :
   $$\limsup_{n \to \infty} \int f_n d\mu \le \int f d\mu \le \liminf_{n \to \infty} \int f_n d\mu$$
   Puisque l'on a toujours $\liminf \le \limsup$, toutes ces valeurs doivent être strictement égales.
   La limite des intégrales existe donc et coïncide avec l'intégrale de la limite. La convergence $L^1$ s'obtient en appliquant ce même raisonnement à la suite $|f_n - f|$ dominée par $2g$.

\section{Applications en Physique, Logique et AI}

- \textbf{Théorie des Probabilités :} Le TCD justifie le passage à la limite sous le signe d'espérance mathématique. Si une suite de variables aléatoires $X_n$ converge p.s. vers $X$, et si elles sont dominées par une variable $Y$ d'espérance finie, alors $\mathbb{E}[X_n] \to \mathbb{E}[X]$.
- \textbf{Apprentissage Automatique (Machine Learning) :}
    - \textbf{Convergence du Gradient Stochastique (SGD) :} Pour prouver que le gradient stochastique calculé sur des mini-lots converge vers le gradient réel de la fonction de perte théorique, le TCD garantit que la dérivation (qui est une limite) peut traverser l'espérance (qui est une intégrale), à condition que la variance des gradients soit dominée.
    - \textbf{Auto-encodeurs Variationnels (VAE) :} La maximisation de l'ELBO nécessite de calculer des gradients d'espérances. Le \textit{Reparameterization Trick} déplace le gradient sous l'intégrale, une opération formellement sécurisée par le TCD.
- \textbf{Mécanique Statistique :} Le calcul des fonctions de partition implique des sommes et intégrales sur des états énergétiques. Le TCD valide les dérivations thermodynamiques (calcul d'entropie, chaleur spécifique) par rapport à la température.


\newpage
\part{Exercices d'Application et de Concours}
\subsection*{{Exercice 1 : Application directe : exponentielle \quad $\bigstar\bigstar\star\star\star$}}

\textbf{Énoncé :}
Calculer la limite suivante :
$$ \lim_{n \to \infty} \int_0^{1} \left(1 - \frac{x}{n}\right)^n e^{x/2} dx $$

\textbf{Correction :}
1. Soit $f_n(x) = (1 - x/n)^n e^{x/2} \mathbf{1}_{[0, 1]}(x)$. Pour $x \in [0, 1]$ fixé et $n$ assez grand, $\ln(1 - x/n) \sim -x/n$, donc $n \ln(1 - x/n) \to -x$. Ainsi, $\lim_{n \to \infty} f_n(x) = e^{-x} e^{x/2} = e^{-x/2}$. La fonction limite est $f(x) = e^{-x/2}$.
2. Cherchons une domination. Pour $t \ge 0$, on sait que $1 - t \le e^{-t}$. Donc $(1 - x/n)^n \le e^{-x}$ pour $0 \le x \le n$.
Ainsi, $|f_n(x)| \le e^{-x} e^{x/2} = e^{-x/2}$ sur $[0, 1]$.
3. La fonction $g(x) = e^{-x/2} \mathbf{1}_{[0, 1]}(x)$ est continue donc mesurable, et intégrable sur $[0, 1]$ (car bornée sur un compact).
4. D'après le Théorème de Convergence Dominée (TCD), on peut intervertir limite et intégrale :
$$ \lim_{n \to \infty} \int_0^1 f_n(x) dx = \int_0^1 e^{-x/2} dx = \left[ -2 e^{-x/2} \right]_0^1 = 2(1 - e^{-1/2}). $$


\subsection*{Exercice 2 : Convergence vers zéro \quad $\bigstar\bigstar\star\star\star$}

\textbf{Énoncé :}
Calculer $\lim_{n \to \infty} \int_0^\infty \frac{\sin(x^n)}{1 + x^2} dx$.

\textbf{Correction :}
Découpons l'intégrale en deux : sur $[0, 1]$ et sur $]1, +\infty[$.
Soit $f_n(x) = \frac{\sin(x^n)}{1 + x^2} \mathbf{1}_{]0, \infty[}(x)$.
1. Pour $x \in [0, 1[$, $x^n \to 0$, donc $\sin(x^n) \to 0$. La limite simple est 0.
Pour $x > 1$, $x^n \to \infty$, la fonction n'a pas de limite simple partout. Cependant, le TCD classique s'applique sur des domaines où la limite existe. Considérons plutôt $\int_0^1 \frac{\sin(nx)}{1 + nx^2} dx$.
1. Soit $f_n(x) = \frac{\sin(nx)}{1 + nx^2}$ sur $]0, 1]$. Pour $x > 0$ fixé, $\lim_{n \to \infty} f_n(x) = 0$.
2. Domination : $|f_n(x)| \le \frac{1}{1 + nx^2} \le \frac{1}{nx^2}$.
Pour une domination indépendante de $n$, notons que $|f_n(x)| \le \frac{n x}{1 + nx^2}$. Sur $\mathbb{R}^+$, $\frac{nx}{1+nx^2} \le \frac{nx}{2\sqrt{n}x} = \frac{\sqrt{n}}{2}$, pas une bonne domination.
Cependant, $|f_n(x)| \le \frac{1}{2x}$ qui n'est pas intégrable en 0.
Un TCD ne s'applique pas directement. Utilisons plutôt $x = u/n$.
$\int_0^n \frac{\sin(u)}{1 + u^2/n} \frac{du}{n}$. Ceci converge vers 0 par changement de variable.
Correction par changement de variable :
Soit $I_n = \int_0^1 \frac{\sin(nx)}{1 + nx^2} dx$. Posons $t = \sqrt{n} x$, $dx = dt/\sqrt{n}$.
$|I_n| \le \int_0^1 \frac{1}{1 + nx^2} dx = \int_0^{\sqrt{n}} \frac{1}{1 + t^2} \frac{dt}{\sqrt{n}} \le \frac{1}{\sqrt{n}} \int_0^\infty \frac{dt}{1+t^2} = \frac{\pi}{2\sqrt{n}}$.
Ainsi, $\lim_{n \to \infty} I_n = 0$. Le TCD n'était pas la meilleure arme ici.


\subsection*{{Exercice 3 : Limite d'une intégrale avec fonction puissance \quad $\bigstar\bigstar\bigstar\star\star$}}

\textbf{Énoncé :}
Montrer que $\lim_{n \to \infty} \int_0^1 \frac{n x^{n-1}}{1 + x} dx = \frac{1}{2}$.

\textbf{Correction :}
Si on utilise le TCD directement, la fonction $f_n(x) = \frac{n x^{n-1}}{1 + x}$ converge simplement vers 0 sur $[0, 1[$. Mais $\int_0^1 f_n(x) dx$ ne tend pas vers 0 ! Il n'y a donc pas de domination possible.
Procédons autrement pour appliquer le TCD.
Par intégration par parties :
$$ \int_0^1 \frac{n x^{n-1}}{1 + x} dx = \left[ \frac{x^n}{1+x} \right]_0^1 - \int_0^1 x^n \left(-\frac{1}{(1+x)^2}\right) dx = \frac{1}{2} + \int_0^1 \frac{x^n}{(1+x)^2} dx $$
Maintenant, appliquons le TCD sur l'intégrale restante.
Soit $h_n(x) = \frac{x^n}{(1+x)^2}$.
1. Convergence simple : pour $x \in [0, 1[$, $\lim_{n \to \infty} h_n(x) = 0$.
2. Domination : $|h_n(x)| \le \frac{1}{(1+x)^2} \le 1$. La constante 1 est intégrable sur $[0, 1]$.
3. Par TCD, $\lim_{n \to \infty} \int_0^1 h_n(x) dx = \int_0^1 0 dx = 0$.
Par conséquent, la limite de l'intégrale initiale est $\frac{1}{2} + 0 = \frac{1}{2}$.


\subsection*{{Exercice 4 : Intégrabilité et continuité de la transformée de Fourier \quad $\bigstar\bigstar\bigstar\star\star$}}

\textbf{Énoncé :}
Soit $f \in \mathcal{L}^1(\mathbb{R})$. On définit sa transformée de Fourier par $\hat{f}(\xi) = \int_{\mathbb{R}} f(x) e^{-i \xi x} dx$.
Montrer à l'aide du TCD que $\hat{f}$ est une fonction continue sur $\mathbb{R}$.

\textbf{Correction :}
Soit $\xi_0 \in \mathbb{R}$ et $(\xi_n)$ une suite convergeant vers $\xi_0$.
Considérons les fonctions $g_n(x) = f(x) e^{-i \xi_n x}$.
1. \textbf{Convergence simple :} Comme l'exponentielle complexe est continue, pour tout $x \in \mathbb{R}$, on a :
   $\lim_{n \to \infty} g_n(x) = f(x) e^{-i \xi_0 x}$.
2. \textbf{Domination :} On a $|g_n(x)| = |f(x)| |e^{-i \xi_n x}| = |f(x)| \times 1 = |f(x)|$.
   Comme $f \in \mathcal{L}^1(\mathbb{R})$, la fonction $|f|$ est intégrable sur $\mathbb{R}$ et indépendante de $n$. Elle constitue une fonction dominatrice parfaite.
3. \textbf{Conclusion :} D'après le TCD, on peut intervertir limite et intégrale :
   $\lim_{n \to \infty} \hat{f}(\xi_n) = \lim_{n \to \infty} \int_{\mathbb{R}} g_n(x) dx = \int_{\mathbb{R}} \lim_{n \to \infty} g_n(x) dx = \int_{\mathbb{R}} f(x) e^{-i \xi_0 x} dx = \hat{f}(\xi_0)$.
   Ceci prouve que $\hat{f}$ est séquentiellement continue en tout point, et donc continue.


\subsection*{{Exercice 5 : Un contre-exemple (masse fuyante) \quad $\bigstar\bigstar\star\star\star$}}

\textbf{Énoncé :}
Soit $f_n(x) = n \mathbf{1}_{]0, 1/n[}(x)$ sur l'espace $[0, 1]$ muni de la mesure de Lebesgue.
1. Calculer $\lim_{n \to \infty} \int_0^1 f_n(x) dx$.
2. Calculer $\int_0^1 \lim_{n \to \infty} f_n(x) dx$.
3. Pourquoi le Théorème de Convergence Dominée ne s'applique-t-il pas ici ?

\textbf{Correction :}
1. Pour chaque $n$, $\int_0^1 f_n(x) dx = n \times \frac{1}{n} = 1$. Donc la limite est 1.
2. Pour $x > 0$ fixé, dès que $n > 1/x$, $f_n(x) = 0$. Donc la limite simple est $f(x) = 0$. L'intégrale de la limite est 0.
3. On a interversion illégitime ($1 \neq 0$). Le TCD ne s'applique pas car il n'existe pas de fonction dominatrice $g$ intégrable telle que $f_n \le g$ pour tout $n$.
En effet, si une telle $g$ existait, on aurait $g(x) \ge \sup_n f_n(x) = 1/x$. Or $x \mapsto 1/x$ n'est pas intégrable sur $[0, 1]$.


\subsection*{{Exercice 6 : Dérivation sous le signe somme (TCD continu) \quad $\bigstar\bigstar\bigstar\bigstar\star$}}

\textbf{Énoncé :}
Soit $F(t) = \int_0^\infty e^{-tx} \frac{\sin x}{x} dx$ pour $t > 0$.
Montrer que $F$ est dérivable sur $]0, +\infty[$ et calculer $F'(t)$.

\textbf{Correction :}
Posons $f(x, t) = e^{-tx} \frac{\sin x}{x}$.
1. $t \mapsto f(x, t)$ est dérivable et $\frac{\partial f}{\partial t}(x, t) = -x e^{-tx} \frac{\sin x}{x} = -e^{-tx} \sin x$.
2. Fixons $a > 0$. Pour tout $t \ge a$ et $x > 0$, on a la domination :
   $\left| \frac{\partial f}{\partial t}(x, t) \right| = e^{-tx} |\sin x| \le e^{-ax}$.
3. La fonction $g(x) = e^{-ax}$ est intégrable sur $[0, +\infty[$.
Par le TCD (version paramétrique, théorème de Leibniz), $F$ est dérivable sur $[a, +\infty[$ et :
$F'(t) = \int_0^\infty -e^{-tx} \sin x dx$.
Calculons cette intégrale en utilisant la partie imaginaire de $\int_0^\infty e^{(-t+i)x} dx$ :
$\int_0^\infty e^{(-t+i)x} dx = \left[ \frac{e^{(-t+i)x}}{-t+i} \right]_0^\infty = \frac{1}{t-i} = \frac{t+i}{t^2+1}$.
La partie imaginaire est $\frac{1}{t^2+1}$. Donc $F'(t) = -\frac{1}{t^2+1}$.


\subsection*{{Exercice 7 : Série de fonctions et interversion \quad $\bigstar\bigstar\bigstar\star\star$}}

\textbf{Énoncé :}
Montrer que $\int_0^1 \frac{x \ln(x)}{1-x} dx = -\sum_{n=1}^{\infty} \frac{1}{(n+1)^2}$.

\textbf{Correction :}
1. Pour $x \in ]0, 1[$, on peut développer $\frac{1}{1-x}$ en série géométrique : $\frac{1}{1-x} = \sum_{n=0}^{\infty} x^n$.
2. L'intégrande devient $\sum_{n=0}^{\infty} x^{n+1} \ln(x)$. Les termes sont tous du même signe (négatif). On peut appliquer le Théorème de Convergence Monotone ou intégrer $-x^{n+1}\ln(x)$ qui est positif.
On peut aussi utiliser le TCD sur les sommes partielles. Les fonctions sont de signe constant, donc le théorème d'intégration terme à terme s'applique (équivalent au TCM).
3. Calculons $\int_0^1 x^{n+1} \ln(x) dx$ par intégration par parties :
   $u(x) = \ln(x)$, $v'(x) = x^{n+1} \implies u'(x) = 1/x$, $v(x) = \frac{x^{n+2}}{n+2}$.
   $\int_0^1 x^{n+1} \ln(x) dx = \left[ \frac{x^{n+2}}{n+2} \ln(x) \right]_0^1 - \int_0^1 \frac{x^{n+1}}{n+2} dx = 0 - \frac{1}{(n+2)^2}$.
4. On a donc $\int_0^1 \frac{x \ln(x)}{1-x} dx = \sum_{n=0}^{\infty} \frac{-1}{(n+2)^2} = -\sum_{k=2}^{\infty} \frac{1}{k^2}$.
Si l'énoncé demande à partir de $n=1$, c'est avec $k=n+1$ donc $-\sum_{n=1}^{\infty} \frac{1}{(n+1)^2}$.


\subsection*{{Exercice 8 : Limite d'une suite définie par intégrale \quad $\bigstar\bigstar\bigstar\star\star$}}

\textbf{Énoncé :}
Calculer $\lim_{n \to \infty} \int_0^n \left(1 - \frac{x}{n}\right)^n \ln(x) dx$.

\textbf{Correction :}
On réécrit l'intégrale sur $]0, +\infty[$ avec une fonction indicatrice :
$I_n = \int_0^\infty f_n(x) dx$ avec $f_n(x) = \left(1 - \frac{x}{n}\right)^n \ln(x) \mathbf{1}_{[0, n]}(x)$.
1. \textbf{Convergence simple :} On sait que $\lim_{n \to \infty} \left(1 - \frac{x}{n}\right)^n = e^{-x}$. Ainsi, la limite simple est $f(x) = e^{-x} \ln(x)$.
2. \textbf{Domination :} On utilise l'inégalité $1 - u \le e^{-u}$ pour $u \in [0, 1]$.
   Donc $\left(1 - \frac{x}{n}\right)^n \le e^{-x}$ pour $0 \le x \le n$.
   On en déduit que $|f_n(x)| \le e^{-x} |\ln(x)| \mathbf{1}_{[0, n]}(x) \le e^{-x} |\ln(x)|$.
3. La fonction $g(x) = e^{-x} |\ln(x)|$ est-elle intégrable sur $]0, +\infty[$ ?
   En 0, $g(x) \sim |\ln(x)|$, dont l'intégrale impropre converge (primitive $x \ln(x) - x$).
   En $+\infty$, $x^2 g(x) = x^2 e^{-x} |\ln(x)| \to 0$, donc $g(x) = o(1/x^2)$, et l'intégrale converge.
4. Par conséquent, par le TCD, la limite de l'intégrale est l'intégrale de la limite :
   $\lim_{n \to \infty} I_n = \int_0^\infty e^{-x} \ln(x) dx$. (Cette intégrale vaut $-\gamma$, où $\gamma$ est la constante d'Euler-Mascheroni).


\subsection*{{Exercice 9 : Domination délicate et TCD \quad $\bigstar\bigstar\bigstar\bigstar\star$}}

\textbf{Énoncé :}
Montrer que la fonction $G(t) = \int_0^\infty e^{-tx^2} \cos(x) dx$ est de classe $C^1$ sur $]0, +\infty[$.

\textbf{Correction :}
Posons $h(x, t) = e^{-tx^2} \cos(x)$.
1. $t \mapsto h(x, t)$ est dérivable et $\frac{\partial h}{\partial t}(x, t) = -x^2 e^{-tx^2} \cos(x)$.
2. Cherchons une domination locale. Soit $[a, b] \subset ]0, +\infty[$ avec $a > 0$.
   Pour $t \in [a, b]$, on a :
   $\left| \frac{\partial h}{\partial t}(x, t) \right| = x^2 e^{-tx^2} |\cos(x)| \le x^2 e^{-ax^2}$.
3. La fonction $g(x) = x^2 e^{-ax^2}$ est intégrable sur $[0, +\infty[$ (décroissance exponentielle très rapide l'emportant sur le polynôme).
4. Le TCD pour la dérivation implique que $G$ est dérivable sur tout $[a, b]$, donc sur $]0, +\infty[$, et que $G'$ est continue car $\frac{\partial h}{\partial t}$ est dominée par $g$ et continue par rapport à $t$. $G$ est donc de classe $C^1$.


\subsection*{{Exercice 10 : Lemme de Scheffé (Corollaire du TCD) \quad $\bigstar\bigstar\bigstar\bigstar\bigstar$}}

\textbf{Énoncé :}
Soit $(f_n)$ une suite de fonctions de $\mathcal{L}^1(\mu)$ et $f \in \mathcal{L}^1(\mu)$.
Supposons que :
1. $f_n \to f$ presque partout.
2. $\int |f_n| d\mu \to \int |f| d\mu$.
Montrer que $f_n$ converge vers $f$ dans $L^1$, c'est-à-dire $\lim_{n \to \infty} \int |f_n - f| d\mu = 0$.

\textbf{Correction :}
On ne peut pas appliquer le TCD directement car on n'a pas de fonction dominatrice explicite.
1. Considérons la fonction $h_n = |f_n| + |f| - |f_n - f|$.
   Par l'inégalité triangulaire, $|f_n - f| \le |f_n| + |f|$, donc $h_n \ge 0$.
2. Convergence simple : Puisque $f_n \to f$ p.p., $|f_n| \to |f|$ et $|f_n - f| \to 0$ p.p.
   Ainsi, $h_n \to |f| + |f| - 0 = 2|f|$ p.p.
3. On applique le lemme de Fatou à la suite de fonctions positives $h_n$ :
   $\int \liminf h_n d\mu \le \liminf \int h_n d\mu$.
   $\int 2|f| d\mu \le \liminf \int (|f_n| + |f| - |f_n - f|) d\mu$.
4. Par linéarité, le côté droit vaut :
   $\liminf \left( \int |f_n| d\mu + \int |f| d\mu - \int |f_n - f| d\mu \right)$.
   Puisque $\int |f_n| d\mu \to \int |f| d\mu$ par hypothèse, on a :
   $\liminf (\dots) = 2\int |f| d\mu + \liminf \left( - \int |f_n - f| d\mu \right)$.
   $2\int |f| d\mu \le 2\int |f| d\mu - \limsup \int |f_n - f| d\mu$.
5. En soustrayant la quantité finie $2\int |f| d\mu$, on obtient :
   $0 \le - \limsup \int |f_n - f| d\mu$, ce qui implique $\limsup \int |f_n - f| d\mu \le 0$.
   Comme l'intégrale est positive, la limite existe et vaut $0$.




\newpage
\part{Travaux Pratiques et Simulations Algorithmiques}
\subsection*{TP 1 : Illustration du TCD sur une intégrale simple}

\begin{lstlisting}
# TP 1 : Illustration numerique du Theoreme de Convergence Dominee
import math

# On etudie l'integrale de f_n(x) = (1 - x/n)^n * exp(x/2) sur [0, 1]
# La limite simple est f(x) = exp(-x/2)

def fn(x, n):
    if x >= n: return 0.0
    return (1.0 - x/n)**n * math.exp(x/2.0)

def f_limite(x):
    return math.exp(-x/2.0)

def integrale(func, x_min=0, x_max=1, num_points=1000):
    dx = (x_max - x_min) / num_points
    total = 0.0
    for i in range(num_points):
        x = x_min + i * dx + dx/2
        total += func(x) * dx
    return total

print("Approximation des integrales I_n :")
for n in [1, 2, 5, 10, 50, 100]:
    val = integrale(lambda x: fn(x, n))
    print(f"n = {n:3} : I_n = {val:.6f}")

val_limite = integrale(f_limite)
print(f"Integrale de la limite f : {val_limite:.6f}")
val_exacte = 2.0 * (1.0 - math.exp(-0.5))
print(f"Valeur theorique exacte  : {val_exacte:.6f}")
assert abs(val_limite - val_exacte) < 1e-4, "L'approximation numerique est erronee"
\end{lstlisting}

\subsection*{TP 2 : Défaut de Domination : Contre-Exemple}

\begin{lstlisting}
# TP 2 : Le phenomene de la bosse glissante (masse fuyante)
# La convergence simple est verifiee, mais l'absence de domination empeche le TCD

# Fonction indicatrice ponderee f_n(x) = n sur ]0, 1/n[
def fn(x, n):
    if 0 < x < 1.0 / n:
        return float(n)
    return 0.0

# La limite simple est 0 partout
def f_limite(x):
    return 0.0

def integrale(func, n, num_points=5000):
    dx = 1.0 / num_points
    total = 0.0
    for i in range(num_points):
        x = i * dx + dx/2
        total += func(x, n) * dx
    return total

print("Analyse de la bosse glissante :")
for n in [10, 50, 100, 500]:
    val = integrale(fn, n)
    print(f"n = {n:3} : I_n = {val:.6f}")

print("La limite des integrales est clairement 1.0")
print("Cependant, l'integrale de la limite f(x)=0 est 0.0")
print("Conclusion : sans domination, limite de l'integrale != integrale de la limite")
\end{lstlisting}

\subsection*{TP 3 : Domination Intégrable : Le Couvercle}

\begin{lstlisting}
# TP 3 : Visualisation du role de la fonction dominatrice g(x)

import math

def fn(x, n):
    # f_n(x) = sin(nx) / (1 + x^2)
    return math.sin(n * x) / (1.0 + x * x)

def g_dom(x):
    # La domination naturelle est 1 / (1 + x^2)
    return 1.0 / (1.0 + x * x)

def abs_fn(x, n):
    return abs(fn(x, n))

print("Verification empirique de la domination : |f_n(x)| <= g(x)")
pts = [0.1, 0.5, 1.0, 2.0, 5.0]

for n in [1, 5, 10]:
    print(f" n = {n}:")
    for x in pts:
        val_f = abs_fn(x, n)
        val_g = g_dom(x)
        print(f"  x = {x:.1f} -> |f_n| = {val_f:.4f}  |  g = {val_g:.4f}")
        assert val_f <= val_g + 1e-9, "Domination violee!"

print(" fonction g(x) forme un couvercle absolu sur toutes les courbes f_n.")
\end{lstlisting}

\subsection*{TP 4 : Dérivation sous le signe intégral (TCD paramétrique)}

\begin{lstlisting}
# TP 4 : Validation numerique de l'application du TCD a la derivation
# On considere F(t) = int_{0^{inf} exp(-tx) * (sin x / x) dx}

import math

def integrand(x, t):
    if abs(x) < 1e-7: return 1.0
    return math.exp(-t * x) * (math.sin(x) / x)

def integrand_derivee(x, t):
    return -math.exp(-t * x) * math.sin(x)

def calc_integral(func, t, x_max=100.0, num_points=5000):
    dx = x_max / num_points
    total = 0.0
    for i in range(num_points):
        x = i * dx + dx/2
        total += func(x, t) * dx
    return total

t = 1.0
dt = 0.001

# Taux d'accroissement
F_t = calc_integral(integrand, t)
F_tdt = calc_integral(integrand, t + dt)
deriv_num = (F_tdt - F_t) / dt

# Integrale de la derivee
deriv_theo = calc_integral(integrand_derivee, t)
val_exacte = -1.0 / (t*t + 1.0) # cf exercice 6

print(f"Derivee par taux d'accroissement (differences finies) : {deriv_num:.6f}")
print(f"Integrale de la derivee partielle (Leibniz/TCD)       : {deriv_theo:.6f}")
print(f"Valeur exacte formelle (-1 / (t^2+1))                 : {val_exacte:.6f}")
\end{lstlisting}

\subsection*{TP 5 : Convergence en norme L1}

\begin{lstlisting}
# TP 5 : Simulation de la convergence L1 via le TCD
# Si f_n converge p.p vers f et est dominee, alors ||f_n - f||_1 tend vers 0.

import math

# f_n(x) = (nx)/(1+n^2 x^2) sur [0, 1]
# lim simple f(x) = 0 pour x > 0. En x=0, f_n(0)=0.
def fn(x, n):
    return (n * x) / (1.0 + n*n * x*x)

def f_lim(x):
    return 0.0

def erreur_L1(n, num_points=1000):
    dx = 1.0 / num_points
    err = 0.0
    for i in range(num_points):
        x = i * dx + dx/2
        # |f_n - f| = f_n(x)
        err += abs(fn(x, n) - f_lim(x)) * dx
    return err

print("Erreur L1 : int_0^1 |f_n(x) - f(x)| dx")
for n in [1, 10, 50, 100, 500]:
    err = erreur_L1(n)
    print(f"n = {n:3} : ||f_n - f||_1 = {err:.6f}")

print("'erreur tend vers 0. La fonction est dominee par g(x) = 1/(2x) mais ")
print("attention, 1/(2x) n'est pas integrable en 0. Le TCD direct est delicat,")
print("pourtant la convergence L1 a bien lieu par un calcul direct d'integrale.")
\end{lstlisting}



\end{document}
